\section{Pendahuluan: Konteks dan Urgensi}
\label{sec:pendahuluan}

\subsection{Latar Belakang}

Pemanasan global akibat aktivitas manusia meningkatkan risiko cuaca ekstrem, kekeringan, dan kenaikan muka laut \cite{P-IPCC-2023}. Di Indonesia, sektor ketenagalistrikan masih bertumpu pada PLTU batu bara. Sistem Jawa-Madura-Bali mengoperasikan 57 unit PLTU terhubung jaringan PLN dengan kapasitas 30,8 GW \cite{D-GEM-GCPT}. Emisinya juga membawa polutan udara yang menurut CREA dan IESR menyebabkan sekitar 10.500 kematian dini pada 2022 \cite{P-CREA-IESR-2023}.

Kerangka hukum untuk mengakhiri PLTU sudah tersedia melalui Perpres 112/2022 dan Permen ESDM 10/2025 \cite{R-PERPRES-112-2022,R-PERMEN-ESDM-10-2025}. CIPP JETP menargetkan pensiun dini 1,7 GW hingga 2040 \cite{P-JETP-CIPP-2023,P-CREA-CIPP-2023}. Di sisi energi bersih, pemerintah menetapkan kuota PLTS atap 2 GW hingga 2028, dengan 1,85 GW di Jawa-Madura-Bali \cite{B-CNBC-ROOFTOP-2028}. Realisasinya baru 538 MWp pada Juli 2025 \cite{B-KONTAN-ROOFTOP-2025}. Sementara itu, proyek percontohan pensiun dini Cirebon-1 dibatalkan pada Desember 2025 \cite{B-BISNIS-CIREBON-BATAL-2025,B-PETROMINDO-CIREBON-BATAL}.

Laporan Kelompok 17 menjawab situasi ini dengan memilih integrasi PLTS atap dan pensiun dini PLTU sebagai solusi \cite{L-LAPORAN-K17}. Solusi tersebut memuat empat mekanisme, yaitu penghentian PLTU bertahap, substitusi kapasitas melalui PLTS atap, penyebaran beban investasi melalui insentif dan kewajiban pemasangan, serta jaring pengaman sosial dan pelatihan ulang. Pelaksanaannya diatur dalam roadmap empat fase. Dokumen kajian ini menurunkan solusi tersebut menjadi program kebijakan yang dapat dihitung, dijalankan, dan dipantau.

\subsection{Rumusan Masalah}

\begin{enumerate}
  \item Mengapa PLTS atap dan pensiun dini PLTU belum berjalan bersamaan, dan celah kebijakan apa yang menyebabkannya?
  \item Bagaimana rancangan program yang memasangkan keduanya sehingga dapat diterima PLN dan negara?
  \item Seberapa layak program tersebut dari sisi fiskal, keandalan, kesehatan, dan sosial, serta bagaimana tahapan pelaksanaannya?
\end{enumerate}

\subsection{Tujuan dan Sasaran Pengguna}

Kajian ini bertujuan untuk (1) memetakan celah kebijakan yang memisahkan PLTS atap dan pensiun dini, (2) merumuskan program pasangan kuota PLTS atap dan pensiun dini beserta aturan pelaksanaannya, dan (3) menilai kelayakan serta menyusun langkah yang dapat langsung ditindaklanjuti.

Dokumen ini ditujukan kepada Kementerian ESDM sebagai penyusun regulasi dan kuota, Kementerian Keuangan sebagai pengelola insentif dan pendanaan, PT PLN (Persero) sebagai pelaksana kajian pensiun dini dan operator sistem, serta Satuan Tugas Transisi Energi yang diusulkan laporan kelompok sebagai koordinator lintas kementerian.

\subsection{Ruang Lingkup}

Kajian difokuskan pada sistem Jawa-Bali karena menampung sebagian besar PLTU dan penjualan listrik nasional. PLTU yang dikaji adalah unit operasi terhubung jaringan PLN, tidak termasuk PLTU captive industri. Horizon simulasi adalah 2026-2040, sejalan dengan target CIPP JETP. PLTS atap diperlakukan sebagai pengganti energi, bukan pengganti daya mampu saat beban puncak malam hari.

\subsection{Metode Kajian}

Kajian ini memakai analisis kebijakan berbasis bukti dalam empat langkah. Pertama, studi data sekunder dari Global Energy Monitor, Statistik PLN 2025, Global Solar Atlas, CREA dan IESR, serta regulasi dan pemberitaan 2022-2026. Kedua, analisis celah kebijakan dengan membandingkan regulasi yang berlaku terhadap kebutuhan solusi terpilih. Ketiga, simulasi tahunan 2026-2040 dengan model PowerShift yang menghitung neraca PLN, APBN, dana luar, keandalan, emisi, dan kematian dini terhadap baseline phase down di akhir kontrak. Keempat, penyusunan instrumen pelaksanaan berupa SOP, rancangan peraturan, dan lembar penilaian kandidat.

Angka dalam kajian ini mengutamakan sumber primer atau sumber yang dapat diperiksa ulang. Beberapa angka laporan kelompok dimutakhirkan dengan data yang lebih baru atau lebih rinci, misalnya kuota PLTS atap per wilayah dan kompensasi pensiun dini yang dikalibrasi ke nilai refinancing Cirebon-1. Daftar pemutakhiran tersebut dicantumkan pada Lampiran~\ref{lamp:pemutakhiran}.
